\section{Capítulo~\ref{sec:sleep}. El sueño}

Los modos del sueño, las repeticiones y la reducción de las sinapsis están medidos con registros de neuronas, microdiálisis y microscopía electrónica; el horario del sueño y el vaivén lento se describen con los modelos del libro, y la función del sueño es sobre todo materia de hipótesis.

{\footnotesize
\setlength{\tabcolsep}{3pt}\renewcommand{\arraystretch}{1.15}
\begin{longtable}{>{\raggedright}p{0.5cm}>{\raggedright}p{4.0cm}>{\raggedright}p{5.6cm}>{\raggedright}p{2.3cm}>{\raggedright\arraybackslash}p{1.7cm}}
\toprule
N.º & Afirmación & Experimento y qué se midió & Fuente & Estado \\
\midrule\endfirsthead
\toprule
N.º & Afirmación & Experimento y qué se midió & Fuente & Estado \\
\midrule\endhead
1 & Durante el sueño las neuronas de la corteza no callan; cambia el modo de trabajo & Registros largos de neuronas corticales en ratas: en el sueño lento la frecuencia sigue siendo distinta de cero, y los periodos de actividad alternan con periodos de silencio; tras una vigilia larga la frecuencia es mayor en todos los estados, y tras el sueño, menor & \cite{Vyazovskiy2009} & medido \\
2 & \nt{ACET} es alto de día y en el sueño REM, y bajo en el sueño lento (tabla~\ref{tab:sleepmodes}) & Microdiálisis en la corteza de gatos en libre movimiento: en la vigilia la liberación de acetilcolina es un $100$--$175\%$ mayor que en el sueño lento, y en el sueño REM está más o menos al nivel de la vigilia tranquila & \cite{Marrosu1995} & medido \\
3 & \nt{NORA} y \nt{SERO} se apagan en el sueño lento y casi callan en el sueño REM & Registro de neuronas \nt{NORA} individuales en ratas y de neuronas \nt{SERO} en gatos por fases del sueño: la frecuencia es máxima en la vigilia, menor en el sueño lento y casi nula en el sueño REM & \cite{AstonJonesBloom1981,McGintyHarper1976} & medido \\
4 & En el sueño REM los músculos están desconectados por inhibición a través de \nt{GABA} y \nt{GLYC} & En ratas se midió la actividad de las neuronas \nt{ACET} del músculo masetero y se aplicaron bloqueantes directamente sobre ellas: la parálisis solo desaparece si se bloquean a la vez los receptores \nt{GABA} de ambos tipos y los receptores \nt{GLYC} & \cite{BrooksPeever2012} & medido \\
5 & La presión de sueño $S$ es un condensador con dos umbrales~\eqref{eq:sleeppressure}, $\tau_r=18.2$~h, $\tau_d=4.2$~h & Modelo de dos procesos; las constantes de tiempo se obtienen de cómo la potencia de las ondas lentas del EEG crece con la vigilia y cae durante el sueño, y la forma de los umbrales, del momento del despertar espontáneo & \cite{Borbely1982,Daan1984} & teoría \\
6 & La máquina llega por sí sola a $16.1$~h de vigilia y $8.0$~h de sueño (fig.~\ref{fig:twoprocess}) & Cálculo según~\eqref{eq:sleeppressure} con umbrales oscilantes, script \texttt{models/sleep\_models.py}: $16.1$~h de vigilia y $7.9$--$8.0$~h de sueño & --- & modelo \\
7 & Tras $40$~h sin dormir, $S=0.90$, pero el sueño dura solo $8.8$~h: la deuda se devuelve con profundidad, no con duración & Modelo: \texttt{models/sleep\_models.py} da $S=0.90$ y $8.8$~h. Experimento: tras $40.5$~h sin dormir, en ocho personas la potencia de las ondas lentas del EEG en la primera noche de recuperación es significativamente mayor de lo habitual & \cite{Borbely1981} & modelo \\
8 & Físicamente, $S$ es la adenosina & Microdiálisis en gatos: la adenosina extracelular en una de las regiones que controlan la vigilia aumenta durante la vigilia, sigue aumentando con la privación de sueño y cae durante el sueño. Que $S$ sea solo la adenosina no está demostrado & \cite{PorkkaHeiskanen1997,Dunwiddie2001} & hipótesis \\
9 & En el sueño lento la corteza oscila: actividad durante unos cientos de ms y silencio, menos de una vez por segundo ($<1$~Hz; bajo anestesia, sobre todo $0.3$--$0.4$~Hz) & Registro intracelular en gatos anestesiados: despolarización de $0.8$--$1.5$~s y una larga hiperpolarización, ritmo $<1$~Hz, sobre todo $0.3$--$0.4$~Hz; la misma alternancia de actividad y silencio se ve en el sueño natural (fila~1) & \cite{Steriade1993,Vyazovskiy2009} & medido \\
10 & El vaivén es un grupo con realimentación y fatiga~\eqref{eq:updown}; con $b=5$, periodo de $1.1$~s (valor del modelo) y actividad alrededor del $35\%$ del tiempo; con $b=1$, actividad uniforme (fig.~\ref{fig:updown}) & Cálculo, script \texttt{models/sleep\_models.py}: periodo de $1.11$~s, fracción de tiempo activo $0.35$ (promedio sobre un número entero de periodos); con $b=1$ la fracción activa es $1.0$. El periodo del modelo es más corto que el medido (fila~9) & --- & modelo \\
11 & \nt{ACET} apaga el vaivén y pasa la corteza a un trabajo uniforme & La estimulación de un núcleo de neuronas \nt{ACET} en el gato bloqueó el ritmo lento en el $79\%$ de las neuronas corticales y las pasó a trabajo uniforme, sobre todo porque desaparecieron las hiperpolarizaciones largas. La acetilcolina suprime las corrientes lentas de potasio que producen la fatiga & \cite{SteriadeAmzica1993,McCormick1992} & medido \\
12 & Las ráfagas de ritmo de $7$--$15$~Hz las genera el lazo \nt{GLUT}--\nt{GABA} entre la corteza y un nodo de relevo & En rodajas del nodo de relevo visual del hurón, las ráfagas las producen las ráfagas de rebote de las células de relevo ante la inhibición de un núcleo \nt{GABA} vecino, al que ellas mismas excitan & \cite{vonKrosigk1993,SteriadeMcCormickSejnowski1993} & medido \\
13 & Las repeticiones del día van en avance rápido: en el hipocampo, unas $20$ veces más rápido; en la corteza, $6$--$7$ veces & En el sueño lento, las secuencias repetidas de neuronas del hipocampo van comprimidas en el tiempo. Tras correr por una pista, las secuencias de neuronas de lugar se repitieron en el mismo orden en ráfagas breves de $\sim100$~ms, comprimidas unas $20$ veces; en la corteza, la compresión es de $6$--$7$ veces & \cite{Nadasdy1999,LeeWilson2002,Euston2007} & medido \\
14 & Reforzar la coordinación de las repeticiones con la corteza mejora la memoria (relación causal) & En ratas, tras el aprendizaje, una estimulación ligada a las repeticiones reforzó su acoplamiento con las ondas lentas y las ráfagas de ritmo de la corteza: al día siguiente la memoria era alta, y en los controles, al nivel del azar & \cite{Maingret2016} & medido \\
15 & La función de las repeticiones es el aprendizaje intercalado de la memoria lenta & Teoría de los dos sistemas de aprendizaje y cálculos con redes artificiales: el aprendizaje secuencial estropea lo antiguo, el intercalado no. No hay ningún experimento directo en el cerebro de que las repeticiones sirvan precisamente para eso & \cite{McClelland1995} & hipótesis \\
16 & Tras el sueño, las sinapsis se reducen en proporción a su tamaño; las grandes casi no cambian & Microscopía electrónica tridimensional de $6920$ sinapsis corticales de ratón: el área de contacto tras el sueño es $\sim18\%$ menor que tras la vigilia; la reducción es proporcional al tamaño, y las sinapsis grandes y estables no cambian & \cite{deVivo2017} & medido \\
17 & La tarea principal del sueño es la renormalización de los pesos & Hipótesis de la homeostasis sináptica; las filas 1 y 16 concuerdan con ella, pero no demuestran que sea la función principal & \cite{TononiCirelli2014} & hipótesis \\
18 & El sueño REM es una prueba en banco del modelo del mundo & El modo (tabla~\ref{tab:sleepmodes}) está medido; que la red ejecute el modelo del mundo es una suposición teórica, sin experimento & \cite{Hobson2009} & hipótesis \\
19 & Limpieza del cerebro durante el sueño: cuestión abierta & Un experimento: durante el sueño, el espacio intercelular es un $60\%$ más ancho y la limpieza es más rápida. Otro, con un método de medida distinto: durante el sueño y bajo anestesia, la limpieza es notablemente más lenta. Los resultados se contradicen & \cite{Xie2013,Miao2024} & hipótesis \\
20 & Sueño local: zonas de la corteza de un animal despierto caen brevemente en silencio & En ratas tras una vigilia larga, las neuronas de una zona aislada de la corteza callan brevemente, como en el sueño lento, con una onda lenta en el EEG local; en esos momentos la rata se equivoca más en la tarea & \cite{Vyazovskiy2011} & medido \\
\bottomrule
\end{longtable}}
